\subsection{实指数函数的展开}

由于 \(\mathrm{D}^n (e^x) = e^x\) ，按 \(n\) 阶泰勒展开，\(e^x\) 为

\[
e ^ {x} = 1 + \frac {x}{1 !} + \frac {x ^ {2}}{2 !} + \dots + \frac {x ^ {n}}{n !} + \int_ {0} ^ {x} \frac {(x - t) ^ {n}}{n !} e ^ {t} d t.\tag{1}
\]

这个公式的余项 \(>0\) （当 \(x > 0\) ），且有 \((-1)^{n + 1}\) 的符号（当 \(x < 0\) ）；此外，中值不等式表明

\[
\frac {x ^ {n + 1}}{(n + 1) !} <   \int_ {0} ^ {x} \frac {(x - t) ^ {n}}{n !} e ^ {t} d t <   \frac {x ^ {n + 1} e ^ {x}}{(n + 1) !} \quad \text { for } x > 0\tag{2}
\]

\[
\frac {\left| x ^ {n + 1} \right| e ^ {x}}{(n + 1) !} <   \left| \int_ {0} ^ {x} \frac {(x - t) ^ {n}}{n !} e ^ {t} d t \right| <   \frac {\left| x ^ {n + 1} \right|}{(n + 1) !} \quad \text { for } x <   0\tag{3}
\]

现在我们知道，序列 \((x^{n}/n!)\) 当 n 无限增大时以 0 为极限（对一切 \(x \geqslant 0\) ；Gen.~Top., IV, p.~365）；于是，在 (1) 中固定 x 而让 n 无限增大，由 (2) 与 (3) 就得出

\[
e ^ {x} = \sum_ {n = 0} ^ {\infty} \frac {x ^ {n}}{n !}\tag{4}
\]

并且右端的级数在 \(\mathbf{R}\) 的每个紧区间上绝对且一致收敛。特别地，有公式

\[
e = 1 + \frac {1}{1 !} + \frac {1}{2 !} + \dots + \frac {1}{n !} + \dots\tag{5}
\]

这个公式使我们可以计算出与数 \(e\) 任意接近的有理近似值；我们得到

\[
e = 2. 7 1 8   2 8 1   8 2 8 \dots
\]

精确到 \(1/10^{9}\) 。此外，公式 (5) 证明 e 是无理数 \(^{2}\) （Gen.~Top., IV, p.~375）。

注。由于公式 (1) 的余项 \textgreater0 （当 x \textgreater{} 0 时），故对 x \textgreater{} 0 有

\[
e ^ {x} > 1 + \frac {x}{1 !} + \frac {x ^ {2}}{2 !} + \dots + \frac {x ^ {n + 1}}{(n + 1) !}
\]

因而更有

\[
e ^ {x} > \frac {x ^ {n + 1}}{(n + 1) !}
\]

对每个整数 \(n\) 成立：由此推出，\(e^x / x^n\) 以 \(+\infty\) 为极限（当 \(x\) 无限增大时；对每个整数 \(n\) ）：我们将在第 V 章中用另一种方法再次得到这一结果（V, p.~231）。

